\subsection{Tensor Train Language Model}
\label{sec: TTLM}
\subsubsection{Tensor-Train Decomposition}
Suppose the sequence of indices of words in the text $X$ is  $w_1, w_2, \cdots, w_N$, where  $w_i \in \{1, 2, \cdots, |V|\}$ and its corresponding weight in $\mathcal{A}$ is denoted as $\mathcal{A}_{w_1 w_2\cdots w_N}$. We use TT decomposition to  represent $\mathcal{A}_{w_1 w_2\cdots w_N}$ in the TT format \citep{oseledets2011tensor} as follows: 


\begin{align} 
    \label{eq: TT_index}
    \mathcal{A}_{w_1w_2...w_N} 
    & = \underbrace{\tG^{(1)}_{:, w_1}}_{1\times R_1} \underbrace{\tG^{(2)}_{:, w_2, :}}_{R_1 \times R_2} \cdots \underbrace{\tG^{(N)}_{:, w_N}}_{R_{N-1} \times 1}  \\
    &= \sum_{\alpha_1, \cdots, \alpha_{N-1}}   \etG^{(1)}_{w_1\alpha_1} \etG^{(2)}_{\alpha_1w_2\alpha_2} \cdots \etG^{(N)}_{\alpha_{N-1}w_N}  \nonumber
\end{align}

where the tensors  $\tG^{(t)} \in \mathbb{R}^{R_{t-1} \times |V| \times R_t}$ ($t =1,...,d$, $R_0 = R_N = 1$ by definition) are called \textit{TT cores}, and  $R_k$ for $k = 1,\cdots,N$ are called \textit{TT ranks}. 

Despite its site-dependent TT cores $\tG^{(t)}$ potentially giving it more expressiveness for language modeling, this property currently generates unnecessary obstacles to its applicability, like the choice of $R_t$. Here we follow the convention of considering a special class of TT decompositions \cite{khrulkov2018expressive, miller2021tensor}, i.e. supposing all the intermediate TT cores are equal to each other $\tG = \tG^{(2)},\ldots,\tG^{(N-1)} \in \mathbb{R}^{R \times |V| \times R}$ and $\tG^{(1)} = \tG^{(N)} \in \mathbb{R}^{|V| \times R}$ in Eq. \ref{eq: TT_index}. 

\subsubsection{Definition of TTLM}
%
\begin{figure*}[t]
    \centering
    % \vspace{-15pt}
    \includegraphics[scale=0.5]{figure/TTLM.pdf}
    \caption{a) Tensor Train Language Model based on Eq. \ref{eq: ttlm_general}. b) TT core of TTLM-Tiny. c) TT core of TTLM-Large. The dashed line in the square represents $\mathcal{A}, \Phi(X)$, or $\tG$. Note that the only difference between TTLM-Large and TTLM-Tiny is whether to use tensor $\tW^{eh}$.}
    \label{fig: TTLM}
\end{figure*}

We define \textit{Tensor Train Language Model} (TTLM) as: 
\begin{equation}
\begin{split}
    p(X) = \sum_{i_1, \cdots, i_N = 1}^{|V|} \sum_{\alpha_1, \cdots, \alpha_{N-1} = 1}^R f(x^{(1)})_{i_1}\etG^{(1)}_{i_1\alpha_1} \cdots \\ f(x^{(N)})_{i_N}\etG^{(N)}_{\alpha_{N-1}i_N}  \label{eq: ttlm_general}
\end{split}
\end{equation}

where each $\vf(x^{(t)})$ is a one-hot vector having $w_t = 1$ for at most one $t$, and has zeros elsewhere. The tensor diagram notation of TTLM is shown in Fig. \ref{fig: TTLM}a. Note that  Eq. \ref{eq: ttlm_general} can compute the elements
of $\mathcal{A}$ in the low-dimensional space as Eq. \ref{eq: TT_index} does. This can be observed  if we represent the elements of  $\tG^{(t)}_{:, w_t, :}$ in Eq.\ \ref{eq: TT_index} as:
%
\begin{align}
    \label{eq: single G}
    \etG^{(t)}_{\alpha_{t-1}w_t\alpha_{t}}
    & = \sum_{{i = 1}}^{|V|} f(x^{(t)})_{i}\etG^{(t)}_{\alpha_{t-1}i\alpha_{t}} 
\end{align}
% 
Since  $\mathcal{A}_{w_1w_2...w_N}$  here equals to  $p(X)$, we can derive Eq. \ref{eq: ttlm_general} by inserting Eq. \ref{eq: single G} into Eq.\ \ref{eq: TT_index}.

The critical difference between TTLM defined by Eq. \ref{eq: ttlm_general} and  Eq. \ref{eq: TT_index}  is that TTLM has combined the weights $\mathcal{A}$ and the input data $\Phi(X)$  together, indicating its potential to be used for language modeling tasks.

\subsubsection{Recursive Probability Computation.}
\label{sec:recursive probability}
We recursively unfold the calculation of TTLM in Eq. \ref{eq: ttlm_general} and find that $\tG$ has two sources of ``input'': the information from the previous recursive unfolding, and the input data $\vf(x^{(t)})$ (see Eq. \ref{eq: TTLM_unfold} for a detailed version). 
From this perspective, $\tG$ acts as a bilinear map $\tG: \mathbb{R}^{|V|} \times \mathbb{R}^R \rightarrow \mathbb{R}^R $, and we can regard the information in the previous step as a hidden state $\vh^{(t)}_\text{TTLM}$, given by:
%
\begin{align}
    \label{eq: ttlm cell full}
    \vh^{(t)}_{\textrm{TTLM}}
    & = \vf(x^{(t)})^T\tG  \vh^{(t-1)}_{\textrm{TTLM}} 
\end{align}
%
where $\vf(x^{(t)})$, $\tG$, and  $\vh^{(t-1)}_{\textrm{TTLM}}$ are contracted together (we permute the indices of $\tG$ from $\mathbb{R}^{R \times |V| \times R}$ to $\mathbb{R}^{|V| \times R \times R}$ which does not change the number of indices). 

Utilizing this recursive property, we here provide further details about computing $p(X)$ by TTLM in practice. In language modeling, $p(X)$ is often decomposed using the chain rule \citep{bahl1983maximum} as follows:
%
\begin{align}
    p(X) = \prod_{t=1}^N p(x^{(t)}|x^{(1:t-1)}) \nonumber
\end{align}
%
where $x^{(1:t-1)}$ denotes the text $[x^{(1)}, x^{(2)}, \cdots ,x^{(t-1)}]$.  At time $t$, the output prediction of a model, $\vy^{(t)}\in \mathbb{V}$, is a probability distribution of word $x^{(t)}$ given $x^{(1:t-1)}$. 

In TTLM, we define $\vy^{(t)}$ as follows: 
%
\begin{align}
\label{eq: probability_TT}
    \vy^{(t)} = \psi \left(\tG^{(t)} \vh^{(t-1)}_{\textrm{TTLM}}\right)
\end{align}
%
where $\tG^{(t)}\in \mathbb{R}^{|V| \times R}$ is the last TT core in TT format at time $t$. $\psi$ is any function that ensures that $\vy^{(t)}$ is non-negative and that the conditional probabilities sum to $1$. For the use of TTLM as a component in a larger architecture, $\psi$ can be chosen as a constant scaling function to preserve linearity; for stand-alone use of TTLMs, $\psi$ can be chosen to be any appropriate activation function---in the remainder of the paper, we shall use the $\softmax$ function \cite{bridle1990probabilistic}. Fig. \ref{fig: probability_ttlm} provides an example of a recursive calculation of conditional probability.
%
 \begin{figure}[t]
    \centering
\includegraphics[width=0.40\textwidth]{figure/ConditionalProbabilityTTLM.pdf}
    \caption{Recursive calculation of conditional probability in TTLM. Here we provide an example that given the text $x^{(1:3)}$, $\vy^{(4)}=\psi(\tG^{(4)}\vh_{\textrm{TTLM}}^{(3)})$ where $\vy^{(4)}$ is the probability distribution of word $x^{(4)}$. }
    \label{fig: probability_ttlm}
    % \vspace{-4mm}
\end{figure}
%


We can derive  the definition of $\vy^{(t)}$ in high-dimensional space, if we substitute $\vh^{(t-1)}_{\textrm{TTLM}}$ in  Eq. \ref{eq: probability_TT}  by Eq. \ref{eq: ttlm_general} and Eq. \ref{eq: ttlm cell full}:
%
\begin{align} \small
      \vy^{(t)} &= \psi \left( \sum_{i_1, \cdots, i_{t-1}} \sum_{\alpha_1,\cdots, \alpha_{t-1}} f(x^{(1)})_{i_1}\etG^{(1)}_{i_1\alpha_1} \cdots \tG^{(t)}_{\alpha_{t-1}}  \right) \label{eq: probability_original_1} \\
      % &= \psi \left(\langle \mathcal{A}^{(1:t)}), \bigotimes\limits^{t-1}_{i=1}  \vf(x^{(i)}) \rangle_{t-1} \right) \label{eq: probability_original_2} \\
    &=  \psi \left(\langle \mathcal{A}^{(1:t)}, \Phi(X^{(1: t-1)}) \rangle_{t-1} \right)  \label{eq: probability_original_3}
\end{align}
%
where $\mathcal{A}^{(1:t)}\in \mathbb{V}^{\otimes t}$,  $\Phi(X^{(1: t-1)}) = \bigotimes\limits^{t-1}_{i=1}  \vf(x^{(i)}) \in \mathbb{V}^{\otimes t-1}$ and $\langle \cdot \rangle_{t-1}$ denotes the "generalized inner product" defined in Sec. \ref{sec:preliminaries}. Note that Eq. \ref{eq: probability_original_1} is the low-dimensional form of Eq. \ref{eq: probability_original_3}, similarly to the relationship between Eq. \ref{eq: ttlm_general} and Eq. \ref{eq: general definition}.




By these definitions, there are some interesting properties of TTLM. (1) We can use \emph{teacher forcing} \citep{marcus1998rethinking} to learn parameters of TT cores. (2)  The hidden-to-output  tensor $\tG^{(t)}$ is defined to be the same as the input-to-hidden tensor $\tG^{(1)}$. (3) $\tG$ and $\tG^{(t)}$ have no parameters in common. We provide a detailed explanation of the  relationship between different TT cores in Appendix \ref{sec: relationship}. 



